
% ─────────────────────────────────────────────────────────────────────────────
\chapter{Introduction}
\label{chap:introduction}
% ─────────────────────────────────────────────────────────────────────────────

The integration of \gls{ai}-based systems into \gls{df} workflows addresses
the need to examine growing volumes of digital data and to identify content
that may be relevant to an investigation.

When automated support is used to classify, organize, or retrieve images,
however, its operational value depends on the stability and interpretability of
the outputs produced under realistic conditions. Degraded, distributionally
different, or intentionally manipulated inputs may prevent relevant content
from receiving appropriate review priority or may increase the amount of
irrelevant material presented to the analyst.

Against this background, the present thesis evaluates the operational
robustness of transparent proxy models and selected black-box software
configurations used for the automated classification or semantic retrieval of
images depicting firearms and related visual content.

The evaluation compares behavior across clean binary inputs, a separate
\gls{ood} branch, adversarial conditions, and anti-forensic transformations.
Its purpose is not to validate automation as an autonomous evidentiary
decision-making mechanism, but to examine the conditions under which it may
provide controlled and traceable support for forensic triage.

Human oversight, source-to-output traceability, explicit configuration
management, and cautious interpretation of automated results remain central
principles throughout the study
\cite{costantini2019digital,sanna2024preliminary,sanna2025pixelperfect}.

% ─────────────────────────────────────────────────────────────────────────────
\section[Digitalization, Threats, and the Central Role of Digital Forensics]
{Digitalization, Threats, and the Central Role of \glsentrylong{df}}
% ─────────────────────────────────────────────────────────────────────────────

The progressive digitalization of contemporary society has transformed how
information is produced, stored, transmitted, and analyzed. Cloud
infrastructures, distributed services, interconnected mobile devices, social
platforms, and heterogeneous information systems generate continuously
increasing volumes of digital data.

These data play a central role in personal, professional, institutional, and
judicial activities. The same technological evolution has expanded the attack
surface of information systems, increased the opportunities available to
hostile actors, and multiplied the points at which potentially relevant digital
information may be altered, concealed, destroyed, or otherwise compromised
\cite{garfinkel2010digital,casey2011digital}.

Alongside the opportunities created by technological innovation, increasingly
heterogeneous threats have emerged. These include data exfiltration and
manipulation, alteration of multimedia files, selective deletion of traces,
anonymization, encryption, information hiding, and other practices intended to
obstruct the reliable reconstruction of events.

Within this context, \gls{df} assumes a central role as the discipline
concerned with identifying, acquiring, preserving, examining, analyzing, and
presenting digital evidence according to principles of methodological rigor,
integrity, and verifiability
\cite{palmer2001road,reith2002examining,nist2006guide,iso27037}.

According to a definition widely cited in the literature, \gls{df} applies
scientific methods and technical procedures to digital data in order to support
its use within investigative or judicial processes
\cite{palmer2001road}.

Each phase must be conducted in a traceable, repeatable, and technically
defensible manner. The objective is not merely to obtain an analytical result,
but to preserve the possibility of reconstructing how that result was produced,
which data were examined, which operations were performed, and which
limitations affected the analysis.

The forensic process is commonly described through phases including
identification, acquisition, preservation, examination, analysis, and
presentation. Although their precise organization may vary, all of these phases
share requirements concerning integrity, transparency, documentation, and
verifiability
\cite{reith2002examining,nist2006guide}.

These requirements become particularly important when evidence is distributed
across multiple devices, stored in virtualized or remote environments, available
only in volatile form, or exposed to intentional manipulation
\cite{nist2006guide,iso27037}.

The growing incidence of cybercrime and the pervasive role of digital
technologies in the commission, facilitation, and documentation of criminal
conduct have consolidated the function of \gls{df} as a connection between
technical analysis and investigative fact-finding.

Within this framework, the quality of the adopted procedures and the reliability
of the tools used may directly affect the ability to reconstruct events,
identify relevant relationships, and preserve the analytical value of digital
evidence
\cite{costantini2019digital,clough2015principles,europol2019iocta}.

% ─────────────────────────────────────────────────────────────────────────────
\section[Automation of Forensic Analysis and the Role of AI]
{Automation of Forensic Analysis and the Role of \glsentryshort{ai}}
% ─────────────────────────────────────────────────────────────────────────────

The quantitative and qualitative expansion of the data subject to investigation
has made exclusively manual forensic analysis increasingly difficult to
sustain. Traditional methods remain essential for technical control,
verification, and final interpretation, but they may be difficult to scale when
analysts must examine large collections of images, videos, logs, conversations,
system files, and heterogeneous artifacts originating from multiple sources.

Under these conditions, automation can support the organization,
prioritization, and preliminary examination of evidence collections
\cite{garfinkel2010digital,costantini2019digital}.

The integration of \gls{ai} techniques into \gls{df} tools has introduced
additional possibilities for supporting investigative analysis. In particular,
\gls{ml} and \gls{dl} methods may assist with tasks such as:

\begin{itemize}
\item semantic classification of images and videos;
\item recognition of recurring patterns;
\item grouping of visually or semantically similar content;
\item analysis of logs and structured artifacts;
\item detection of anomalies;
\item retrieval and prioritization of material potentially relevant to an
investigation.
\end{itemize}

Automated triage may also affect the conditions under which analysts review
sensitive material. Studies involving professionals engaged in the examination
of child sexual abuse material have documented the potential psychological
impact of prolonged exposure and the operational value of filtering or safer
presentation technologies intended to reduce unnecessary exposure
\cite{sanchez2019practitioner}.

Although this use case differs from the
\texttt{weapon}/\texttt{non\_weapon} task evaluated in the present thesis, it
illustrates that the reliability of automation may affect not only processing
efficiency but also analyst workload and the organization of human review
\cite{costantini2019digital,sanna2024preliminary}.

The application of \gls{ai} to \gls{df} does not merely accelerate existing
procedures. It may alter how investigative information is filtered, ordered,
and prioritized. Systems integrated into forensic software may suggest semantic
categories, group multimedia content, identify apparent anomalies, or retrieve
items associated with specific analyst-defined concepts.

Such functionality can influence which items receive early attention and which
remain outside the initial review set. \gls{ai}-based modules are therefore
increasingly incorporated into professional software to support classification,
sorting, semantic retrieval, and preliminary evidence selection
\cite{sanna2025cybercrime,sanna2025pixelperfect}.

The benefits of automation cannot, however, be assessed exclusively in terms of
speed or workload reduction. A system that assigns categories or returns
semantically related images does not simply accelerate examination. Its outputs
may affect the order in which evidence is reviewed and the resources allocated
to particular items.

A false negative may prevent relevant content from being prioritized, whereas a
false positive may increase analytical noise and the amount of material
requiring manual verification. The operational consequences therefore depend
on both the frequency and the direction of the errors.

The introduction of \gls{ai} into forensic workflows consequently requires
critical evaluation of:

\begin{itemize}
\item nominal task performance;
\item behavior under unexpected or manipulated inputs;
\item false-negative and false-positive distributions;
\item system version and configuration;
\item output semantics and normalization rules;
\item traceability from the source item to the automated result;
\item the continued role of expert verification.
\end{itemize}

\gls{xai} methods may provide supplementary diagnostic information for
transparent models by representing the attribution assigned to selected input
features. Such visualizations do not automatically establish prediction
correctness, causal relevance, semantic understanding, or forensic reliability.

Explainability must therefore remain subordinate to quantitative validation and
expert interpretation
\cite{Adadi2018,Doshi-Velez2017,sanna2025cybercrime}.

% ─────────────────────────────────────────────────────────────────────────────
\section[Robustness of AI-Based Systems and Adversarial and Anti-Forensic Threats]
{Robustness of \glsentryshort{ai}-Based Systems and Adversarial and
Anti-Forensic Threats}
% ─────────────────────────────────────────────────────────────────────────────

The use of \gls{ai}-based systems in high-stakes domains has highlighted a
fundamental limitation: strong performance under standard conditions does not
necessarily imply stable behavior in the presence of altered, degraded,
distributionally different, or hostile inputs.

This issue is particularly relevant to \gls{df}, where automated
image-analysis systems may process content that has undergone repeated
compression, resizing, color alteration, degradation, platform-specific
processing, or intentional manipulation. They may also encounter images that
differ substantially from the data used during model development or task
definition.

Intentional alteration of digital evidence predates the introduction of
\gls{ai}. Anti-forensic techniques broadly include strategies intended to
obstruct the identification, acquisition, preservation, examination, or
interpretation of potentially relevant data.

Examples include selective deletion of traces, metadata manipulation,
encryption, file concealment, obfuscation, and information hiding.
Steganography, for example, may conceal information within apparently ordinary
digital images, making the hidden content difficult to identify both for human
analysts and for automated tools not specifically designed to detect it
\cite{yaacoub2021digital,sanna2025cybercrime}.

In the present thesis, steganography is referenced as a representative example
of the broader anti-forensic domain. The experimental anti-forensic branch is
more narrowly focused on image-level transformations directly relevant to the
robustness of automated visual analysis.

The literature on \gls{advml} has demonstrated that \gls{ml} and \gls{dl}
models may be vulnerable to inputs deliberately optimized to induce erroneous
predictions. These inputs, commonly referred to as \emph{adversarial examples},
may cause a classifier to fail even when the applied alteration is constrained
or visually limited.

Such findings challenge the assumption that high accuracy on clean data is
sufficient evidence of reliability under realistic or hostile conditions
\cite{goodfellow2015explaining,carlini2017towards,biggio2018wild,nowroozi2021survey}.

Adversarial attacks and anti-forensic techniques must nevertheless remain
conceptually distinct. Adversarial attacks are defined in relation to their
effect on a model or decision process and may require different levels of
knowledge or access, depending on the adopted threat model.

Anti-forensic techniques have a broader objective. They may alter, conceal, or
degrade digital evidence independently of whether an \gls{ai}-based classifier
is present. Their effectiveness is not necessarily defined by a particular
model response.

The two perspectives intersect when a manipulation of digital evidence reduces
the ability of an automated system to identify, retrieve, or prioritize content
of investigative interest. In this situation, a manipulation directed at an
automated classifier may also have an anti-forensic consequence because it
affects triage and subsequent human review
\cite{sanna2024preliminary,sanna2025pixelperfect}.

The threat model is therefore essential to the interpretation of robustness
results. A perturbation generated directly against an accessible model,
empirical transfer of the same artifact to a different architecture, a
score-based attack, and a deterministic model-agnostic transformation represent
different experimental conditions.

Their effects must not be interpreted as equivalent forms of attack success or
as interchangeable evidence of system robustness.

The notion of operational robustness adopted in this thesis refers to the
observable stability of a defined system, version, configuration, task, and
output mapping when evaluated across the frozen input conditions included in
the protocol.

Robustness is not treated as an absolute or version-independent property. A
system may perform strongly on clean inputs while exhibiting:

\begin{itemize}
\item increased missed-target risk under a specific manipulation;
\item additional false-positive review workload;
\item closed-set assignments on out-of-task content;
\item sensitivity to thresholds, semantic-query settings, or output mappings;
\item severe degradation under a direct-target adversarial condition.
\end{itemize}

Operational robustness therefore concerns not only aggregate performance but
also the quality, direction, traceability, and workflow consequences of the
automated support provided to the analyst.

% ─────────────────────────────────────────────────────────────────────────────
\section{Research Problem and Relevance in Realistic Forensic Scenarios}
% ─────────────────────────────────────────────────────────────────────────────

Despite the extensive literature on adversarial attacks and model robustness,
the evaluation of these issues in operationally oriented forensic scenarios
remains comparatively limited. Many studies rely on standardized benchmarks,
general-purpose \textit{computer vision} datasets, or experimental conditions
that do not reproduce the organization and interpretive constraints of a
forensic workflow.

Studies more closely aligned with forensic practice indicate that automated
image analysis must address large evidence collections, heterogeneous sources,
task-specific semantic categories, ambiguous content, and outputs intended to
support analyst review.

Silva et al.\ propose an \gls{ai}-based system for classifying images contained
in criminal evidence, including categories such as firearms, ammunition,
documents, screenshots, and nudity. More recent work by Sanna et al.\ indicates
that automated functions integrated into forensic software may exhibit
operationally relevant limitations when processing manipulated or unexpected
content
\cite{silva2021microservices,sanna2024preliminary,sanna2025pixelperfect}.

In \gls{df}, the research problem is not limited to demonstrating that a
classifier can be induced to fail. It is also necessary to determine how an
error affects the support provided to the analyst.

Relevant questions include whether a manipulation:

\begin{itemize}
\item prevents target content from receiving appropriate review priority;
\item increases the volume of irrelevant material presented for examination;
\item alters the balance between recall and selectivity;
\item produces closed-set assignments on content outside the nominal task;
\item exploits a specific model directly or transfers to other systems;
\item reveals sensitivity to configuration or output-normalization choices.
\end{itemize}

Addressing these questions requires evaluation populations aligned with the
defined operational task. Such populations should include realistic positives,
realistic negatives, manipulated variants, borderline cases, and content that
does not conform to the nominal classification problem.

This requirement is consistent with research concerning the evaluation of
automated visual-content systems. Although it addresses perceptual hashing
rather than firearm classification, the PHASER framework emphasizes that
evaluation datasets should reflect the intended deployment scenario and its
relevant anomalous cases
\cite{mckeown2024phaser}.

The present thesis applies this principle to a task centered on realistic
individual firearms and clearly recognizable firearm-related visual content.
This category is of potential investigative interest, has been considered in
previous research, and permits controlled experimental validation
\cite{silva2021microservices,Hao2018DeepFirearm}.

The supervised task comprises the operational classes
\texttt{weapon} and \texttt{non\_weapon}. A separate \gls{ood} branch contains
borderline, anomalous, degraded, synthetic, semantically adjacent, and otherwise
out-of-task content.

The \gls{ood} branch is not a third supervised class and is not used to train a
three-class classifier. It is analyzed separately to characterize the
assignments produced when closed-set systems process inputs that do not conform
to the nominal binary task.

Adversarial and anti-forensic artifacts are generated exclusively from the
1\,000-image binary subset. This preserves the distinction between:

\begin{itemize}
\item supervised binary correctness;
\item clean-to-perturbed error transitions;
\item separate \gls{ood} assignment behavior.
\end{itemize}

The evaluation design combines two complementary levels of observation. The
first uses transparent proxy models, for which predictions, maximum
predicted-class probabilities, error transitions, checkpoints, and internal
gradients can be inspected under controlled conditions.

The second evaluates selected black-box software configurations through the
outputs observable to the analyst. These include tags, bookmarks, exported
classification labels, and semantic-retrieval membership. No assumption is made
that the proxy models reproduce or explain the proprietary internal models used
by those systems.

The use of the same frozen source population and derived perturbation suite
provides a common experimental basis while preserving the distinction between
transparent model evaluation and operational black-box assessment.

Within this framework, vulnerability has direct workflow implications.
A \texttt{weapon}-to-\texttt{non\_weapon} transition may prevent potentially
relevant content from being prioritized. A
\texttt{non\_weapon}-to-\texttt{weapon} transition may increase manual review
workload. An \gls{ood} weapon assignment is not an ordinary binary false
positive, but it may still route out-of-task content toward a firearm-oriented
review category.

The research problem is therefore located at the intersection of algorithmic
robustness, software configuration, output interpretation, and operational
reliability. The objective is not merely to measure generic performance
degradation, but to determine how automated image-analysis systems behave under
a frozen set of realistic and potentially hostile conditions and how the
resulting outputs may affect forensic triage.

The scientific relevance of the study lies in the definition and application of
an integrated, task-specific evaluation protocol. Its operational relevance
lies in identifying the limitations, error profiles, configuration dependencies,
and traceability requirements that must remain visible when automated tools are
used to support evidence analysis.

The thesis consequently contributes to research on domain-specific evaluation
and comparative protocols for \gls{ai} in \gls{df}, while preserving human
verification as the final interpretive layer
\cite{silva2021microservices,mckeown2024phaser,sanna2024preliminary,sanna2025pixelperfect}.


% ─────────────────────────────────────────────────────────────────────────────
\section{Research Objectives}
% ─────────────────────────────────────────────────────────────────────────────

In light of the critical issues identified in the previous sections, this
thesis aims to evaluate systematically the operational robustness of
transparent proxy models and selected black-box software systems used for the
automated classification or semantic retrieval of images of potential
investigative interest in \gls{df}.

The general objective is to determine whether the behavior observed under clean
reference conditions remains stable when the evaluated systems process
distributionally different, degraded, transformed, or intentionally perturbed
images. In an investigative context, such alterations may reduce the likelihood
that relevant content is prioritized through automated detection or triage
procedures or may increase the amount of irrelevant material requiring human
review.

The experimental task distinguishes two supervised operational classes,
\texttt{weapon} and \texttt{non\_weapon}, from a separate \gls{ood} branch.
The latter contains external, ambiguous, borderline, degraded, synthetic, or
otherwise out-of-task images and is not treated as a third supervised class.
Its purpose is to characterize the closed-set assignments produced when the
input does not conform to the nominal binary task
\cite{hendrycks2017baseline,hendrycks2019benchmarking}.

Adversarial attacks and anti-forensic transformations are also treated as
distinct experimental branches. The former include model-dependent
perturbations designed to affect classifier decisions and the deterministic
model-agnostic \gls{colorshift} condition. The latter represent plausible
image-level modifications that may arise through ordinary media processing or
may be intentionally used to reduce the effectiveness of automated analysis.

Broader anti-forensic behaviors, including the use of general-purpose
generative \gls{ai} tools to facilitate the creation of steganographic code,
remain outside the experimental scope of this thesis. They nevertheless
illustrate how automated technologies may expand the range of techniques
available for concealing information in digital images
\cite{sanna2025cybercrime}.

Within this scope, the work is organized around four main objectives.

The first objective is to define a rigorous, documented, and reproducible
methodological protocol for evaluating operational robustness in a realistic
forensic image-analysis scenario. The protocol includes human-validated dataset
construction, deterministic partitioning, transparent proxy-model evaluation,
controlled perturbation generation, preparation of a blind forensic evaluation
bundle, black-box output normalization, quantitative analysis, qualitative
explainability, and traceability controls.

The second objective is to evaluate heterogeneous systems on the same frozen
source population and derived perturbation suite. The transparent evaluation
level comprises EfficientNet-B0, ResNet18, and a \gls{clip}-based proxy
consisting of a frozen visual encoder and a supervised binary classification
head. The operational black-box level comprises the selected software
configurations evaluated exclusively through outputs observable to the analyst
and through explicitly documented normalization rules.

This parallel design does not assume that the proxy models reproduce or
approximate the proprietary internal components of the selected software
systems. Instead, it provides an inspectable experimental reference and a
common basis for comparing observable operating profiles.

The third objective is to characterize robustness through more than aggregate
accuracy. The analysis considers false negatives, false positives,
clean-to-perturbed error transitions, \gls{ood} weapon assignment rates,
architecture-specific maximum predicted-class probabilities, and
configuration-dependent changes in recall and selectivity.

For the adversarial branch, the evaluation explicitly distinguishes attacks
applied directly to the EfficientNet-B0 source models from their empirical
transfer to ResNet18, the \gls{clip}-based proxy, and the selected black-box
configurations. The deterministic \gls{colorshift} condition remains separate
because it does not rely on gradients, model parameters, checkpoints, or
observable scores.

A complementary qualitative \gls{xai} analysis is conducted on purposively
selected EfficientNet-B0 cases. Its purpose is to support case-level diagnostic
inspection rather than to provide quantitative robustness evidence, causal
feature explanations, or information about the internal behavior of the
black-box systems.

The fourth objective is to relate the experimental findings to the operational
use of automation in \gls{df}. Automated classifications, tags, bookmarks, and
semantic retrieval results are treated as triage indicators rather than as
autonomous evidentiary determinations.

The analysis therefore aims to identify the limitations, error profiles,
configuration dependencies, traceability requirements, and forms of human
verification necessary for the use of these systems to remain compatible with
a \textit{human-in-the-loop} forensic workflow
\cite{costantini2019digital,faqir2023digital,sanna2024preliminary}.

Overall, the research is methodological and operational rather than exclusively
performance-oriented. It examines how interactions among data, models,
perturbations, software configurations, output mappings, and analyst review
affect the support provided by automated image-analysis systems in realistic
and potentially hostile conditions.






% ─────────────────────────────────────────────────────────────────────────────
\subsection{Research Questions}
\label{subsec:introduction-research-questions}
% ─────────────────────────────────────────────────────────────────────────────

The experimental design addresses the following research questions:

\begin{enumerate}
    \item[\textbf{RQ1.}] To what extent do the transparent proxy models maintain
    their clean-baseline performance under the separate \gls{ood}, adversarial,
    and anti-forensic conditions?

    \item[\textbf{RQ2.}] How do direct source-model attacks, cross-architecture
    transfer, and the model-agnostic \gls{colorshift} condition differ in their
    effects on the evaluated proxy models?

    \item[\textbf{RQ3.}] How do the selected black-box software systems respond
    to the same frozen experimental corpus, and how do observable output
    semantics and configuration choices affect the resulting operational
    profiles?

    \item[\textbf{RQ4.}] Which error transitions, robustness degradations,
    \gls{ood} assignment behaviors, and non-operational outputs are most
    relevant to forensic triage and continued human review?
\end{enumerate}


% ─────────────────────────────────────────────────────────────────────────────
\section{Contribution of the Thesis}
% ─────────────────────────────────────────────────────────────────────────────

The main contribution of this thesis lies in the definition and application of
a methodological protocol for evaluating the operational robustness of
transparent proxy models and selected black-box software systems used for
automated image classification or semantic retrieval in \gls{df}.

Unlike evaluations focused exclusively on standard \textit{computer vision}
benchmarks or aggregate performance under clean conditions, this work adopts a
forensic and operational perspective. It examines how automated systems behave
when they process clean binary inputs, a separate \gls{ood} population,
adversarial artifacts, and anti-forensic transformations
\cite{nowroozi2020phd,hendrycks2019benchmarking,sanna2024preliminary,sanna2025pixelperfect}.

The contribution of the thesis is articulated across four interconnected
dimensions.

The first contribution is the formalization of the
\textbf{\gls{fairlab}} protocol. \gls{fairlab} integrates:

\begin{itemize}
\item human-validated dataset construction;
\item deterministic data partitioning;
\item transparent proxy-model evaluation;
\item adversarial and anti-forensic artifact generation;
\item preparation of a blind forensic evaluation bundle;
\item black-box software assessment;
\item explicit output-normalization rules;
\item quantitative error-oriented metrics;
\item qualitative explainability analysis;
\item traceability and integrity controls.
\end{itemize}

\gls{fairlab} is not proposed as a new classifier, defensive mechanism,
standalone forensic application, or general theoretical architecture. It is a
documented procedure through which the behavior of heterogeneous automated
systems can be evaluated under controlled but operationally relevant
conditions.

Its methodological value lies in the coordinated application of the individual
components and in the preservation of the relationships among source images,
derived artifacts, system outputs, normalized decisions, metrics, and reporting
assets.

The second contribution is the construction of a frozen experimental corpus
aligned with the selected forensic image-analysis task. The source population
contains:

\begin{itemize}
\item 500 \texttt{weapon} images;
\item 500 \texttt{non\_weapon} images;
\item 500 images reserved for the separate \gls{ood} branch.
\end{itemize}

The \texttt{weapon} branch contains realistic individual firearms and clearly
recognizable firearm-related visual content. The \texttt{non\_weapon} branch
contains realistic negatives consistent with the binary task. The separate
\gls{ood} branch includes borderline, anomalous, degraded, synthetic,
semantically adjacent, and otherwise out-of-task images.

The \gls{ood} population is not treated as a third supervised class. It is used
to characterize the closed-set assignments produced when an input does not
conform to the nominal \texttt{weapon}/\texttt{non\_weapon} task.

The 1\,000-image binary subset was used to generate five adversarial or
adversarial-style populations and five anti-forensic populations. Together with
the clean binary images and the separate \gls{ood} branch, these derivatives
form an 11\,500-item forensic evaluation bundle:

\begin{itemize}
\item 1\,000 clean binary items;
\item 500 clean \gls{ood} items;
\item 5\,000 adversarial artifacts;
\item 5\,000 anti-forensic artifacts.
\end{itemize}

The 10\,000 manipulated artifacts remain derived versions of the same 1\,000
binary source images. They represent source--condition combinations rather than
independent underlying cases.

Sample selection follows a documented \textit{human-in-the-loop} procedure.
Manual validation is treated as an explicit methodological component rather
than concealed as an automatic ground-truth generation process. The adopted
criteria, reviewer decisions, and associated limitations remain traceable
\cite{sambasivan2021data}.

The third contribution is the parallel evaluation of transparent proxy models
and selected black-box software configurations on the same frozen experimental
basis.

The transparent level comprises EfficientNet-B0, ResNet18, and a
\gls{clip}-based proxy consisting of a frozen visual encoder and a supervised
binary classification head. These models provide access to fold-specific
checkpoints, predicted classes, maximum predicted-class probabilities,
confusion matrices, matched clean-to-perturbed transitions, and internal
gradients.

The black-box level comprises:

\begin{itemize}
\item \gls{magnetaxiomtool}/\gls{magnetai};
\item \gls{griffeye}/\gls{t3kcore};
\item \gls{excirefoto2025};
\item \gls{cellebriteinseyets}.
\end{itemize}

These systems are evaluated exclusively through outputs observable to the
analyst, including tags, bookmarks, exported classification labels, and
semantic-query result membership. Their outputs are normalized through
explicitly documented operational rules without assuming that the systems
implement equivalent internal taxonomies or binary classification tasks.

The proxy models are not intended to reconstruct, approximate, or explain the
proprietary models used by the black-box systems. They provide an inspectable
experimental reference and the controlled generation layer for the
model-dependent adversarial artifacts subsequently evaluated for empirical
transfer.

The protocol therefore distinguishes:

\begin{itemize}
\item direct-target attacks against fold-specific EfficientNet-B0 checkpoints;
\item empirical transfer to ResNet18 and the \gls{clip}-based proxy;
\item empirical transfer to the selected black-box configurations;
\item the deterministic and model-agnostic \gls{colorshift} condition;
\item the separate anti-forensic transformations.
\end{itemize}

This distinction prevents transfer results from being interpreted as direct
robustness evaluations of systems against which no matched attack was
independently optimized.

The fourth contribution is an auditable and error-oriented interpretation of
operational robustness. The evaluation does not reduce system behavior to a
single aggregate score. It considers:

\begin{itemize}
\item false negatives affecting firearm-oriented content;
\item false positives increasing analyst review workload;
\item directional clean-to-perturbed error transitions;
\item closed-set weapon assignments on the separate \gls{ood} branch;
\item direct-target and transfer behavior;
\item configuration-dependent changes in recall and selectivity;
\item residual embedded-metadata sensitivity;
\item the continued need for human verification.
\end{itemize}

Stable identifiers, manifests, fold assignments, model checkpoints,
perturbation records, cryptographic digests, blind-bundle identifiers,
normalized predictions, metric tables, and reporting assets preserve the
relationship between each source item and the corresponding experimental
outcomes.

This traceability supports procedural reproducibility and artifact-level
auditability, including where complete source datasets or proprietary software
exports cannot be redistributed.

A complementary qualitative \gls{xai} analysis is conducted through
\gls{ig} on five purposively selected EfficientNet-B0 cases. The attribution
maps support case-level diagnostic inspection of the selected proxy-model
outputs. They are not treated as quantitative robustness evidence, causal
feature explanations, semantic localization ground truth, or explanations of
the selected black-box systems.

Overall, the contribution of the thesis does not consist in proposing a
universally robust firearm classifier or ranking the evaluated software
products. It consists in defining and applying a traceable procedure for
revealing how heterogeneous automated image-analysis systems behave when the
assumptions of clean, closed-set evaluation no longer hold.

The resulting perspective treats operational robustness as a functional
requirement for the controlled use of automation in \gls{df}. Automated
outputs remain triage indicators whose value depends on task-specific
validation, explicit configuration and mapping rules, source-to-result
traceability, and continued expert review
\cite{costantini2019digital,sanna2024preliminary,sanna2025pixelperfect}.


% ─────────────────────────────────────────────────────────────────────────────
\section{Technical-Forensic Relevance and Operational Reliability}
% ─────────────────────────────────────────────────────────────────────────────

In forensic contexts, the technical performance of an automated tool cannot be
assessed independently of the task for which it is used and of the consequences
that its outputs may have on evidence analysis. The introduction of
\gls{ai}-based systems into investigative workflows therefore raises a question
that is both technical and operational: under which conditions can an automated
output provide useful support to the analyst without reducing the completeness,
traceability, or critical review of the examination
\cite{faqir2023digital,sanna2025cybercrime}?

The relevance of an automated output depends on more than aggregate
classification accuracy. It also depends on the direction of the errors, the
software or model version, the active configuration, the observable output
semantics, the adopted normalization rule, the input population, and the threat
model under which the system was evaluated.

A false negative on firearm-oriented content may prevent a potentially relevant
image from receiving appropriate review priority. A false positive may increase
the amount of irrelevant material requiring manual examination, reduce the
selectivity of the triage process, and divert analyst attention. Neither error
direction can therefore be evaluated independently of the intended workflow and
the relative operational consequences of missed detections and additional
review workload.

The separate \gls{ood} branch introduces a further distinction. An
\gls{ood} image assigned to \texttt{weapon} is not an ordinary supervised false
positive because the input has no binary reference assignment. It may
nevertheless affect the workflow by triggering a flag, changing the priority of
an item, or increasing the amount of material presented to the analyst.

Automated classifications, tags, bookmarks, and semantic retrieval results must
therefore be interpreted as triage indicators rather than as autonomous
evidentiary determinations. They may support the organization and prioritization
of large evidence collections, but they do not independently establish the
nature, relevance, authenticity, or evidentiary weight of an image.

This issue becomes particularly significant when the evaluated systems operate
as black boxes. In such cases, the analyst may observe a tag, bookmark,
classification label, exported percentage, or semantic-query result without
having access to the internal model, decision function, preprocessing pipeline,
or thresholding procedure that produced it.

These observable outputs do not necessarily share the same semantics. They
cannot be assumed to represent equivalent binary classifications or directly
comparable uncertainty measures. Their interpretation must remain associated
with the specific software version, enabled functions, query or category
configuration, export procedure, and normalization rule used in the evaluation.

The same caution applies to the transparent proxy models. Maximum
predicted-class probabilities are architecture-specific and are not calibrated
across models. A high value does not independently establish correctness,
calibrated certainty, operational robustness, or forensic reliability.

The qualitative \gls{xai} component provides a limited form of diagnostic
transparency for selected EfficientNet-B0 outputs. Through \gls{ig}, the
analysis visualizes attribution magnitudes associated with the adopted
predicted-class target, baseline, and integration procedure.

These attribution maps do not identify causal features, certify prediction
correctness, provide semantic localization ground truth, or explain the
proprietary systems evaluated in the black-box branch. They supplement
quantitative error analysis but do not replace robustness testing, convergence
assessment, or expert interpretation
\cite{Adadi2018,Doshi-Velez2017}.

Operational reliability is therefore closely connected to traceability,
verifiability, and critical oversight. An automated result is more defensible
when it can be linked to:

\begin{itemize}
\item the original source item and its integrity record;
\item the model or software version;
\item the active checkpoint or operational configuration;
\item the observable output produced by the system;
\item the normalization rule applied to that output;
\item the experimental or operational condition under which the item was
processed;
\item the subsequent human review and final interpretation.
\end{itemize}

The present thesis is situated within this problem space. It does not assume
that automation is intrinsically reliable and does not attempt to validate an
automated classifier as an autonomous source of technical truth. Instead, it
examines observable behavior under a frozen protocol in order to identify the
conditions under which automated image analysis may provide useful triage
support and those under which particular caution is required.

From a legal and evidentiary perspective, the study remains deliberately
limited. It does not determine the admissibility, probative value, or
evidentiary weight of an automated classification, nor does it certify any of
the evaluated software products. Such determinations would require
case-specific legal analysis, independent validation, chain-of-custody
assessment, disclosure of the applicable technical procedures, and
consideration of the relevant procedural framework.

The technical-forensic relevance of the work instead lies in supporting a
distinction between the controlled instrumental use of \gls{ai} and uncritical
reliance on automated outputs. In \gls{df}, automation remains most defensible
when it operates as a documented, auditable, and reviewable component of a
\textit{human-in-the-loop} workflow
\cite{cpp1988,aiact2024,swgde_guidelines}.



% ─────────────────────────────────────────────────────────────────────────────
\section{Structure of the Thesis}
% ─────────────────────────────────────────────────────────────────────────────

The thesis is organized into seven chapters, progressing from the theoretical
and research context to the methodological protocol, its concrete
implementation, the experimental results, and their operational interpretation.
A separate appendix documents the principal repository and reproducibility
artifacts.

\textbf{Chapter~\ref{chap:introduction}} introduces the research problem,
defines the objectives and Research Questions, presents the main contributions,
and discusses the technical-forensic relevance of evaluating the operational
robustness of transparent proxy models and selected black-box software
configurations.

\textbf{Chapter~\ref{chap:background}} presents the theoretical foundations
related to digital evidence, automated forensic analysis, the use of
\gls{ai} for image classification, \gls{advml}, anti-forensic techniques,
\gls{ood} behavior, and \gls{xai}.

\textbf{Chapter~\ref{chap:state-of-the-art}} critically examines the relevant
literature and identifies the research gap concerning the integrated and
operationally oriented evaluation of transparent models and selected black-box
software systems under clean, distributionally different, adversarial, and
anti-forensic conditions.

\textbf{Chapter~\ref{chap:methodology}} defines the high-level
\gls{fairlab} methodological protocol. It describes the research design,
dataset-construction and semantic-validation principles, the separate binary
and \gls{ood} branches, the parallel proxy-model and black-box evaluation
levels, the threat model, the evaluation metrics, the explainability approach,
and the principles governing traceability, integrity, and reproducibility.

\textbf{Chapter~\ref{chap:implementation}} documents the concrete
implementation of the protocol and the frozen experimental setup. It describes
the data structures, scripts, fold construction, proxy-model configurations,
parameter settings, software versions, perturbation-generation procedures,
forensic evaluation bundle, black-box output-normalization rules,
explainability workflow, and technical validation controls.

\textbf{Chapter~\ref{chap:results}} presents and interprets the experimental
evidence. It reports the clean proxy-model baseline, the separate \gls{ood}
analysis, the anti-forensic and adversarial results, the behavior of the
selected black-box configurations, the embedded-metadata sensitivity check,
the comparative operational analysis, and the qualitative
EfficientNet-B0 cases examined through \gls{xai}.

\textbf{Chapter~\ref{chap:conclusions}} answers the Research Questions,
summarizes the principal contributions and experimental findings, discusses
their operational and forensic implications, defines the limitations of the
study, and identifies corresponding directions for future research.

The separate
\textbf{Appendix~\ref{app:repository-reproducibility-artifacts}} documents the
repository structure, canonical experimental artifacts, official pipeline
sequence, integrity records, provenance information, validation assets, and the
boundaries of controlled reproducibility.




